\scanpage{525}
\mspage{436}
% B01
par \unsure{AQT|SQT} qu'on peut trouver des cardinaux $\pi_0$, \inserted{$\pi_1$} ayant les propriétés suivantes,

% B02
\textcircled{\scriptsize $1^\circ$}) $\pi_0$ est utile pour $E_s$, $\forall s\in\mathrm{Ob}\,C$

% B03
\[ 2^\circ)\quad \pi_1\geq\mathrm{Sup}\left(\underline{\pi_0},\ \underline{\mathrm{card}\,C},\ \mathop{\mathrm{Sup}}\limits_{\lambda\in\mathrm{Fl}\,C}\pi_\lambda,\ \underline{\mathop{\mathrm{Sup}}\limits_{s\in\mathrm{Ob}\,C}\unsure{\mathrm{card\,red}|\mathrm{cardred}}\,E_{s\pi_0}}\right. \]

% B04
\textcircled{\scriptsize $3^\circ$}) $\forall\,(\lambda\inserted{(s\to t)}\in\mathrm{Fl}\,C$ et $x_0\in\mathrm{Ob}\,E_{s\pi_0}$, le foncteur
\[ H_\lambda(x_0)=(x_1\mapsto\mathrm{Hom}_\lambda(x_0,x_1)):E_t\unsure{\relbar|\to}(\mathrm{Ens}) \]
est $\pi_1$-accessible. (NB \deleted{\unsure{Pour ceci, et|Pour ceci, il}} \deleted{\unsure{suffit que}} \inserted{si $f$ est cofibrant}, cel\unsure{a} signifie que $\forall\lambda\in\mathrm{Fl}\,C,\lambda:s\relbar t$, $\underset{\sim}{\boxed{\lambda_*(\unsure{E|S}_{s\pi_0})\subset E_{s\pi_1}}}$.
% B05
Si $f$ est fibrant, il suffit d'exiger que \deleted{$\lambda^*$ soit} $\forall\lambda\in\mathrm{Fl}\,C$, $\lambda^*:E_t\to E_s$ soit $\pi_1$-accessible.) Mais dans ce cas, on \unsure{veut|voit|aura} aussi $\boxed{\lambda^*(E_{t\pi_0})\subset E_{s\pi_1}}$

% B06
\textcircled{\scriptsize $4^\circ$}) les foncteurs $\lambda^*$ (cas $f$ fibrant) resp $\lambda_*$ (cas $f$ cofibrant) sont $\pi_1$-accessibles.

% B07
Soit donc $\pi$ un cardinal tel que
{\csname tagsleft@true\endcsname
\[ \pi=\pi^{\pi_0}\geq\pi_1, \tag*{$(\unsure{*|\alpha})$} \]}
\deleted{je dis que} \inserted{et soit} $E(\pi)$ la sous-catégorie \inserted{\unsure{(str.)|(dit.)|(dir.)}} pleine de $E$ \deleted{\unsure{formée}} définie par
\[ \mathrm{Ob}\,E(\pi)=\bigcup_{s\in\mathrm{Ob}\,C}\mathrm{Ob}\,E_{s\pi}. \]

% B08
Grâce à $(\unsure{*|\alpha})$, et à $1^\circ$\deleted{\unsure{, et}} $4^\circ$ et à
\[ \left|\begin{aligned}
&\pi_1\geq\mathop{\mathrm{Sup}}\limits_s(\mathrm{card\,red}\,E_{s\pi_0})\ \unsure{\times|x}\quad\text{(contenu ds }2^\circ)\\
&\lambda^{*}_{\unsure{*|{}}}(E_{t\pi_0})\subset E_{s\pi_1}\quad\text{resp}\quad\lambda_*(E_{s\pi_0})\subset E_{t\pi_1}\\[-1pt]
&\hspace{16em}\text{(contenu ds }3^\circ)
\end{aligned}\right. \]
on voit que $E(\pi)$ est une sous-catégorie

% archive sheet number: 524
